% !TEX program = xelatex
% Шаблон научной статьи НИТУ «МИСиС»
% Заполните блоки ниже и скомпилируйте: xelatex article.tex

\documentclass{misisarticle}

% --- Классификационные коды (оставьте пустым {}, если не требуются) ---
\grnti{27.41.41, 28.23.29}
\asjc{2608}
\oecd{1.02}

% --- Название ---
\title{Топологическая классификация дефектов кремниевых пластин на основе триангуляции Делоне}

% --- Авторы: \misisauthor{ФИО}{аффилиация}{город, страна}{email} ---
\misisauthor{В.~А.~Попадьина}{кафедра инженерной кибернетики НИТУ «МИСиС»}{Москва, Россия}{m2312071@edu.misis.ru}
\misisauthor{И.~Ю.~Федоров}{кафедра инженерной кибернетики НИТУ «МИСиС»}{Москва, Россия}{}

% --- Аннотация и ключевые слова ---
\abstract{%
  В статье рассматривается оптимизация процесса построения триангуляции Делоне для больших наборов точек, 
  получаемых при сегментации изображений дефектов кремниевых пластин, чтобы обеспечить корректное и быстрое обнаружение 
  и классификацию дефектов. Статья включает в себя сравнительный обзор существующих алгоритмов и инженерных приемов ускорения 
  триангуляции, а также влияние выбора алгоритма на устойчивость последующего топологического анализа и точность классификации дефектов.
}

\keywords{%
  топологический анализ данных, персистентная гомология, триангуляция, дефекты кремниевых пластин, алгоритмы оптимизации
}

\begin{document}

\maketitle

% ============================================================
\section{Введение}

Современная полупроводниковая промышленность предъявляет строгие требования к контролю качества на всех этапах производства 
интегральных микросхем. Непрерывная миниатюризация элементов приводит к увеличению числа и разнообразия дефектов, таких как 
разрывы и перемычки (bridges) в литографических паттернах, что напрямую снижает выход годных изделий. Для оперативного контроля 
применяются сканирующие электронные микроскопы (SEM) и лазерные сканирующие микроскопы (ЛСМ), позволяющие получать высокоточные изображения дефектов.

Однако ручной анализ больших объемов таких изображений неэффективен. Современные методы компьютерного зрения, основанные на глубоком 
обучении, позволяют автоматизировать сегментацию дефектов \cite{lee2024}. Тем не менее, для последующей классификации и анализа причин 
возникновения дефектов требуются методы, способные улавливать структурную целостность объектов. Альтернативой выступает топологический 
анализ данных (ТДА). Использование персистентной гомологии позволяет извлекать устойчивые топологические сигнатуры дефектов, независимые 
от системы координат. В данной работе обосновывается выбор триангуляции Делоне в качестве метода построения симплициального комплекса для ТДА.
% ============================================================
\section{Анализ связанных работ}
Задача автоматизации контроля дефектов в полупроводниковом производстве активно исследуется. В работе \cite{lee2024} предложен метод слабо 
контролируемой сегментации SEM-изображений (WISE-SEM), который позволяет снизить затраты на разметку данных, используя ограничивающие рамки 
вместо попиксельных масок. Авторы отмечают, что для анализа причин дефектов недостаточно только классификации по форме, так как требуется детальное 
выделение границ. Этим обусловлена необходимость перехода к более информативным структурным методам анализа.

В области вычислительной геометрии фундаментальные алгоритмы построения триангуляции Делоне были заложены в работах 
\cite{lee1980, shamos1975}. Ли и Шахтер \cite{lee1980} предложили эффективный алгоритм «разделяй и властвуй» со сложностью \(O(n \log n)\), 
который до сих пор остается одним из самых надежных и предсказуемых методов. Шамо и Хоуи \cite{shamos1975} связали задачу построения 
триангуляции с поиском ближайших соседей, что легло в основу современных вычислительных пакетов.

Среди работ сотрудников кафедры инженерной кибернетики НИТУ «МИСИС» можно выделить исследования, посвященные оценке геометрических 
дефектов методами 3D-реконструкции \cite{demin2026} и определению размеров объектов на изображениях с использованием монокулярной 
оценки глубины \cite{prokhorov2026}. Эти работы демонстрируют важность точного восстановления геометрии объектов для последующего 
анализа. Также в работе научного руководителя \cite{korovin2026} предложена топологически-инспирированная метрика для обнаружения 
потенциальных дефектов в литографии, что напрямую подтверждает перспективность применения методов ТДА в данной предметной области.

\section{Постановка задачи}
Пусть задано множество точек \(S \subset \mathbb{R}^2\), полученных в результате сегментации изображений дефектов. Требуется построить 
триангуляцию \(T(S)\) так, чтобы она удовлетворяла критерию Делоне: ни одна точка из \(S\) не должна находиться внутри описанной 
окружности любого треугольника из \(T(S)\).

Критерии выбора алгоритма построения триангуляции для промышленной системы:
\begin{itemize}
  \item \textbf{Производительность:} гарантированная сложность не выше \(O(|S| \log |S|)\) для обработки сотен тысяч точек.
  \item \textbf{Предсказуемость:} стабильное время работы в худшем случае (worst-case).
  \item \textbf{Устойчивость:} корректная работа с вырожденными случаями (коллинеарные точки, точки на одной окружности).
\end{itemize}

Исходя из этих критериев, необходимо выбрать алгоритм, который обеспечивает баланс между скоростью и надежностью.

\section{Описание методов}

\subsection{Сравнительный анализ алгоритмов триангуляции}

Существует множество алгоритмов построения триангуляции Делоне. В таблице~\ref{tab:algo_compare} приведено их сравнение по основным критериям.

\begin{table}[htbp]
  \caption{Сравнение алгоритмов построения триангуляции Делоне}
  \label{tab:algo_compare}
  \centering
  \begin{tabular}{lccc}
    \toprule
    \textbf{Алгоритм} & \textbf{Сложность} & \textbf{Устойчивость} & \textbf{Основной недостаток} \\
    \midrule
    Полный перебор & \(O(|S|^4)\) & Низкая & Экспоненциально медленный \\
    Алгоритм переворота & \(O(|S|^2)\) & Средняя & Возможны зацикливания \\
    Инкрементальный & \(O(|S|^2)\) & Средняя & Зависимость от порядка точек \\
    Sweep-Hull & \(O(|S| \log |S|)\) & Высокая & Требует сортировки \\
    Разделяй и властвуй & \(O(|S| \log |S|)\) & Высокая & Сложность слияния \\
    \bottomrule
  \end{tabular}
\end{table}

Как видно из таблицы, алгоритм «Разделяй и властвуй» \cite{lee1980} является оптимальным выбором для данной задачи, 
так как обеспечивает детерминированную сложность \(O(|S| \log |S|)\) и высокую численную стабильность.

\subsection{Условие Делоне и построение комплекса}

Триангуляция множества точек \(S\) — это планарный граф, разбивающий выпуклую оболочку \(S\) на треугольники. Условие Делоне 
математически записывается через определитель матрицы, построенной на координатах точек. Для четырех точек \(A, B, C, D\) проверка 
условия того, что точка \(D\) лежит вне описанной окружности треугольника \(ABC\), сводится к вычислению:
\begin{equation}
  \begin{vmatrix}
    x_A & y_A & x_A^2 + y_A^2 & 1 \\
    x_B & y_B & x_B^2 + y_B^2 & 1 \\
    x_C & y_C & x_C^2 + y_C^2 & 1 \\
    x_D & y_D & x_D^2 + y_D^2 & 1 \\
  \end{vmatrix} > 0.
  \label{eq:incircle}
\end{equation}
Если условие~\eqref{eq:incircle} выполняется, то триангуляция локально удовлетворяет критерию Делоне.


\subsection{Связь с персистентной гомологией}

Построенная триангуляция Делоне используется как фильтрация (filtration) для вычисления диаграмм персистентности. Каждому ребру 
и треугольнику присваивается вес (например, длина ребра). В процессе увеличения порогового значения \(\varepsilon\) в комплексе 
появляются и исчезают топологические признаки. Точки на диаграмме персистентности, удаленные от главной диагонали, соответствуют 
значимым дефектам (например, трещинам), что позволяет проводить их классификацию.


% ============================================================
\section{Заключение}

Кратко сформулируйте основные выводы работы.

% ============================================================
% Список литературы — вариант 1: вручную
\begin{thebibliography}{9}

\bibitem{lee2024} 
Lee Y., Kim S. B. Weakly Supervised Image Segmentation for Detecting Defects From Scanning Electron Microscopy Images in Semiconductor // IEEE Access. — 2024. — Vol. 12. — P. 184896–184908. — DOI: 10.1109/ACCESS.2024.3513873.

\bibitem{lee1980} 
Lee D. T., Schachter B. J. Two algorithms for constructing a Delaunay triangulation // International Journal of Computer and Information Sciences. — 1980. — Vol. 9, no. 3. — P. 219–242.

\bibitem{shamos1975} 
Shamos M. I., Hoey D. Closest-point problems // Proceedings of the 16th Annual Symposium on Foundations of Computer Science (SFCS). — 1975. — P. 151–162.

\bibitem{demin2026} 
Демин Н. И., Матвеев В. С. Исследование возможности оценки геометрических дефектов дорожного покрытия методами 3D реконструкции // Искусственный интеллект в промышленных, коммерческих, медицинских и финансовых приложениях : Сборник статей научно-технического семинара кафедры "Инженерной кибернетики". — Москва: НИТУ «МИСиС», 2026. — С. 28-33.

\bibitem{prokhorov2026} 
Прохоров Д. Д. Определение геометрических размеров объектов на изображениях с использованием монокулярной оценки глубины // Искусственный интеллект в промышленных, коммерческих, медицинских и финансовых приложениях : Сборник статей научно-технологического семинара кафедры "Инженерной кибернетики". — Москва: НИТУ «МИСиС», 2026. — С. 98-103.

\bibitem{korovin2026} 
Korovin A. N., Vasilyev A. I., Budennyy S., Fedorov I. Topology-Inspired Metric for Detecting Potential Defects in Lithography // Expert Systems with Applications. — 2026. — DOI: 10.1016/j.eswa.2026.132161.

\end{thebibliography}

\end{document}
